% Template:     Informe LaTeX
% Documento:    Archivo principal - Tarea 2 (Game of Life en CPU, CUDA y OpenCL)
% Codificación: UTF-8

% CREACIÓN DEL DOCUMENTO
\documentclass[
	spanish,
	letterpaper, oneside
]{article}

% INFORMACIÓN DEL DOCUMENTO
\def\documenttitle {Tarea 2: OpenCL y CUDA}
\def\documentsubtitle {Juego de la Vida de Conway en CPU y GPU}
\def\documentsubject {Computación en GPU}

\def\documentauthor {Matías Saavedra}
\def\coursename {Computación en GPU}
\def\coursecode {CC7515-1}

\def\universityname {Universidad de Chile}
\def\universityfaculty {Facultad de Ciencias Físicas y Matemáticas}
\def\universitydepartment {Departamento de Ciencias de la Computación}
\def\universitydepartmentimage {departamentos/dcc}
\def\universitydepartmentimagecfg {height=1.57cm}
\def\universitylocation {Santiago de Chile}

% INTEGRANTES, PROFESORES Y FECHAS
\def\authortable {
	\begin{tabular}{ll}
		Integrante:
		& \begin{tabular}[t]{l}
			Matías Saavedra
		\end{tabular} \\
		Profesora:
		& \begin{tabular}[t]{l}
			Nancy Hitschfeld K.
		\end{tabular} \\
		Auxiliar:
		& \begin{tabular}[t]{l}
			Gustavo Medel
		\end{tabular} \\
		Ayudante:
		& \begin{tabular}[t]{l}
			Catalina Gajardo
		\end{tabular} \\
		& \\
		\multicolumn{2}{l}{Fecha de entrega: \today} \\
		\multicolumn{2}{l}{\universitylocation}
	\end{tabular}
}

% IMPORTACIÓN DEL TEMPLATE
\input{template}

% INICIO DE PÁGINAS
\begin{document}

% PORTADA
\templatePortrait

% CONFIGURACIÓN DE PÁGINA Y ENCABEZADOS
\templatePagecfg

% RESUMEN O ABSTRACT
\begin{abstractd}
	Este informe describe el desarrollo y evaluación de tres implementaciones
	del Juego de la Vida de Conway: una versión secuencial en CPU, una
	versión paralela en CUDA y una versión paralela en OpenCL. Se incluye
	una variación de las reglas originales en la que una célula muerta
	también nace si tiene exactamente $6$ vecinos vivos. Se miden las
	células evaluadas por segundo para distintos tamaños de cuadrícula
	$N \times M$, se comparan dos variaciones de configuración del kernel
	GPU, y se analizan los \emph{speed-up} respecto de la implementación en
	CPU para determinar a partir de qué tamaño de cuadrícula resulta
	conveniente utilizar la GPU.
\end{abstractd}

% TABLA DE CONTENIDOS - ÍNDICE
\templateIndex

% CONFIGURACIONES FINALES
\templateFinalcfg

% ======================= INICIO DEL DOCUMENTO =======================

\section{Introducción}
\label{sec:introduccion}

% (0.5 pt)
%   - Breve descripción del problema.
%   - Resultados esperados.

El Juego de la Vida de Conway es un autómata celular bidimensional en el
que cada celda de una grilla $N \times M$ se encuentra viva o muerta y
evoluciona en pasos discretos según el estado de sus ocho vecinos. En esta
tarea se considera la variación de las reglas estándar dada por:

\begin{itemize}
	\item Una célula nace en un espacio muerto si tiene exactamente $3$ vecinos vivos.
	\item Una célula sobrevive a la próxima generación si tiene $2$ o $3$ vecinos vivos.
	\item Una célula también nace en un espacio muerto si tiene exactamente $6$ vecinos vivos.
\end{itemize}

El objetivo es implementar el algoritmo en CPU (secuencial), CUDA y OpenCL,
medir su rendimiento en células evaluadas por segundo, y analizar el
\emph{speed-up} respecto de la versión secuencial.

\subsection{Resultados esperados}
\label{sec:resultados-esperados}

Se espera que las implementaciones en GPU superen significativamente a la
CPU para grillas grandes, donde el paralelismo de datos se vuelve dominante
frente al costo de transferencias de memoria entre host y device. Para
grillas pequeñas, en cambio, se espera que el costo de lanzamiento del
kernel y de transferencia haga que la CPU sea competitiva o incluso más
rápida.

\section{Implementación}
\label{sec:implementacion}

% (1.0 pt)
%   - Implementación secuencial.
%   - Implementación paralela en CUDA.
%   - Implementación paralela en OpenCL.
%   - Variaciones en configuración.

\subsection{Implementación secuencial en CPU}
\label{sec:impl-cpu}

Describir aquí la estructura del programa secuencial: representación de la
grilla como arreglo plano de \texttt{unsigned char}, doble buffer
(\texttt{current}/\texttt{next}), recorrido por filas y columnas, y conteo
de vecinos con condiciones de borde.

\subsection{Implementación paralela en CUDA}
\label{sec:impl-cuda}

Describir aquí la organización del \emph{grid} y los \emph{blocks}, el
mapeo de un hilo por celda, el uso (o no) de memoria compartida con
\emph{halo}, y la sincronización entre generaciones.

\subsection{Implementación paralela en OpenCL}
\label{sec:impl-opencl}

Describir aquí la creación del contexto, la cola de comandos, la
compilación del kernel y la configuración del \emph{work-group}. Notar las
diferencias con CUDA.

\subsection{Variaciones en configuración}
\label{sec:variaciones}

Se eligieron las siguientes dos opciones a comparar (de las tres
disponibles en el enunciado):

\begin{enumerate}
	\item Tamaños de bloque $32$, $64$, $128$ y $256$.
	\item Uso de memoria local por bloque \emph{vs.} sin memoria local.
\end{enumerate}

% Reemplazar por la elección real si difiere.

\section{Resultados}
\label{sec:resultados}

% (2.0 pt)
%   - Describir experimentos a realizar.
%   - Resultados por tamaño de cuadrícula (CPU/GPU).
%   - Resultados de experimentos para las opciones elegidas.

\subsection{Diseño de experimentos}
\label{sec:experimentos}

Se midió el número de células evaluadas por segundo en grillas cuadradas
de tamaños $N \in \{128, 256, 512, 1024, 2048, 4096\}$, ejecutando un
número fijo de generaciones por configuración y promediando varias
corridas independientes.

\subsection{CPU vs GPU por tamaño de cuadrícula}
\label{sec:res-cpu-gpu}

Insertar aquí la tabla y gráficos comparativos.

\subsection{Resultados de las opciones elegidas}
\label{sec:res-opciones}

Insertar aquí las curvas para cada tamaño de bloque y para la comparación
con/sin memoria local.

\section{Análisis de resultados}
\label{sec:analisis}

% (1.5 pt)
%   - Speed-up comparando versiones paralelas con CPU.
%   - Calcular desde qué tamaño de cuadrícula conviene usar la GPU.

\subsection{Speed-up}
\label{sec:speedup}

Reportar $S(N) = T_{\text{CPU}}(N) / T_{\text{GPU}}(N)$ para cada
configuración.

\subsection{Umbral de conveniencia GPU}
\label{sec:umbral}

Identificar el tamaño de grilla a partir del cual $S(N) > 1$ de manera
sostenida.

\subsection{Perfilamiento (Nsight / Chrono)}
\label{sec:profiling}

Describir hallazgos del \emph{profiling} de la mejor solución CUDA con
Nsight (o Nvprof) y de la mejor solución OpenCL con \texttt{std::chrono}.

\section{Conclusiones}
\label{sec:conclusiones}

% (1.0 pt)

Resumir los principales hallazgos y limitaciones encontradas.

% FIN DEL DOCUMENTO
\end{document}
